\documentclass[12pt,a4paper]{article}

%=========================================================
%                   ENCODAGE ET POLICES
%=========================================================
\usepackage[utf8]{inputenc}
\usepackage[T1]{fontenc}
\usepackage[french]{babel}
\usepackage{lmodern}

%=========================================================
%                     MATHÉMATIQUES
%=========================================================
\usepackage{amsmath,amssymb,amsfonts,mathrsfs}
\usepackage{tkz-tab}
\everymath{\displaystyle}

%=========================================================
%                     MISE EN PAGE
%=========================================================
\usepackage[left=1.5cm,right=1.5cm,top=1.5cm,bottom=1.8cm]{geometry}
\usepackage{multicol}
\usepackage{float}
\usepackage{array,multirow}
\usepackage{enumitem}

%=========================================================
%                GRAPHISMES ET COULEURS
%=========================================================
\usepackage{tikz}
\usetikzlibrary{arrows.meta,calc}
\usepackage{pgfplots}
\pgfplotsset{compat=1.15}

\usepackage[most]{tcolorbox}
\tcbuselibrary{skins,hooks}

%=========================================================
%                  LIENS ET ARRIÈRE-PLAN
%=========================================================
\usepackage[colorlinks=true,linkcolor=blue,urlcolor=blue,citecolor=blue]{hyperref}
\usepackage{eso-pic}

%=========================================================
%                   COULEURS ET STYLES
%=========================================================
\definecolor{myred}{RGB}{180, 40, 40}
\definecolor{myblue}{RGB}{20, 50, 120}
\definecolor{mygreen}{RGB}{30, 120, 40}

%=========================================================
%                       FILIGRANE
%=========================================================
\AddToShipoutPicture{
    \AtTextCenter{
        \makebox[0pt]{
            \rotatebox{45}{
                \textcolor[gray]{0.92}{
                    \fontsize{5cm}{5cm}\selectfont PGB
                }
            }
        }
    }
}

\begin{document}

% En-tête
\begin{center}
    \Large\textbf{\underline{\textcolor{myred}{CHAPITRE 5 : ÉTUDE DE FONCTIONS}}}\\[0.2cm]
    \normalsize\textbf{Prof : M. BA} \hfill \textbf{Classe : Terminale L}\\[0.1cm]
    \textbf{Durée : 6 heures} \hfill \textbf{Année scolaire : 2025 -- 2026}
\end{center}
\hrule
\vspace{0.3cm}

%=========================================================
%      I. PARITÉ ET ÉLÉMENTS DE SYMÉTRIE
%=========================================================

%=========================================================
%      I. PARITÉ ET ÉLÉMENTS DE SYMÉTRIE
%=========================================================

\section*{\color{myred}I. Parité et éléments de symétrie}

Soit $f$ une fonction numérique définie sur un ensemble $D_f$ et $\mathcal{C}_f$ sa courbe représentative dans un repère orthogonal.

\subsection*{\color{myred}1. Parité d'une fonction}

\begin{tcolorbox}[colback=blue!5!white,colframe=myblue,title=\textbf{Définition 1 : Fonction paire}]
On dit qu'une fonction $f$ est \textbf{paire} si :
\begin{enumerate}
    \item Pour tout $x \in D_f$, $-x \in D_f$ ($D_f$ est centré en $0$).
    \item Pour tout $x \in D_f$, $f(-x) = f(x)$.
\end{enumerate}
\textbf{Interprétation géométrique :} L'axe des ordonnées $(Oy)$ est un \textbf{axe de symétrie} pour la courbe $\mathcal{C}_f$.
\end{tcolorbox}

\begin{tcolorbox}[colback=red!5!white,colframe=myred,title=\textbf{Définition 2 : Fonction impaire}]
On dit qu'une fonction $f$ est \textbf{impaire} si :
\begin{enumerate}
    \item Pour tout $x \in D_f$, $-x \in D_f$.
    \item Pour tout $x \in D_f$, $f(-x) = -f(x)$.
\end{enumerate}
\textbf{Interprétation géométrique :} L'origine du repère $O(0,0)$ est un \textbf{centre de symétrie} pour la courbe $\mathcal{C}_f$.
\end{tcolorbox}

\paragraph{Exemple 1 : Étude de la parité}
\begin{itemize}
    \item \textbf{Cas d'une fonction paire :} Soit $f(x) = x^4 - 2x^2 + 1$ définie sur $D_f = \mathbb{R}$.
    \begin{enumerate}
        \item $D_f = \mathbb{R}$ est symétrique par rapport à $0$.
        \item Pour tout $x \in \mathbb{R}$ :
        \[
        f(-x) = (-x)^4 - 2(-x)^2 + 1 = x^4 - 2x^2 + 1 = f(x)
        \]
        Donc $f$ est **paire**, et $(Oy)$ est axe de symétrie de $\mathcal{C}_f$.
    \end{enumerate}

    \item \textbf{Cas d'une fonction impaire :} Soit $g(x) = \dfrac{2x}{x^2 + 1}$ définie sur $D_g = \mathbb{R}$.
    \begin{enumerate}
        \item $D_g = \mathbb{R}$ est symétrique par rapport à $0$.
        \item Pour tout $x \in \mathbb{R}$ :
        \[
        g(-x) = \frac{2(-x)}{(-x)^2 + 1} = \frac{-2x}{x^2 + 1} = -g(x)
        \]
        Donc $g$ est **impaire**, et l'origine $O(0,0)$ est centre de symétrie de $\mathcal{C}_g$.
    \end{enumerate}
\end{itemize}

\vspace{0.3cm}
\hrule
\vspace{0.3cm}
\end{document}